\documentclass[dvisvgm,tikz,border=5pt]{standalone}
\usepackage{amsmath,amssymb}
\usepackage{CJKutf8}
\usetikzlibrary{arrows.meta,positioning,calc}

\definecolor{themebg}{HTML}{30362d}
\definecolor{themefg}{HTML}{edf4e8}
\definecolor{themegray}{HTML}{9ea897}
\definecolor{themecurve}{HTML}{7eb6ff}
\definecolor{themealert}{HTML}{ff7b7b}
\definecolor{themeamber}{HTML}{f5b942}

\pgfdeclarelayer{background}
\pgfdeclarelayer{foreground}
\pgfsetlayers{background,main,foreground}

\begin{document}
\begin{CJK*}{UTF8}{gbsn}
\pagecolor{themebg}
\begin{tikzpicture}[>=Stealth,font=\small,color=themefg,text=themefg]
\tikzset{
  count/.style={rounded corners=3pt,draw=themegray,fill=themefg!3!themebg,minimum width=1.35cm,minimum height=.78cm,align=center},
  bad/.style={count,draw=themealert,fill=themealert!10!themebg},
  good/.style={count,draw=themecurve,fill=themecurve!10!themebg},
  up/.style={->,themegray,line width=.9pt}
}

\node[anchor=west,font=\bfseries] at (0,5.7)
  {层容量题：目标层放多少节点，要把“父节点成本”一层层向上算};
\node[anchor=west,font=\small,text=themegray] at (0,5.15)
  {例：二叉树共 126 个节点，问第 7 层最多能有多少个节点。方框表示各层至少需要的节点数，不是实际树形。};

\node[anchor=west,font=\bfseries,text=themealert] at (0,4.15) {尝试 $x=64$：};
\node[anchor=west,font=\scriptsize,text=themegray] at (2.02,4.83)
  {每条箭头：上一层至少取 $\left\lceil\text{当前层}/2\right\rceil$};
\node[bad] (a7) at (2.15,4.15) {第7层\\64};
\node[count] (a6) at (3.65,4.15) {第6层\\32};
\node[count] (a5) at (5.15,4.15) {第5层\\16};
\node[count] (a4) at (6.65,4.15) {第4层\\8};
\node[count] (a3) at (8.15,4.15) {第3层\\4};
\node[count] (a2) at (9.65,4.15) {第2层\\2};
\node[count] (a1) at (11.15,4.15) {根\\1};
\draw[up] (a7)--(a6);
\draw[up] (a6)--(a5); \draw[up] (a5)--(a4); \draw[up] (a4)--(a3); \draw[up] (a3)--(a2); \draw[up] (a2)--(a1);
\node[anchor=west,font=\small,text=themealert] at (12.05,4.15) {$\sum=127>126$};

\node[anchor=west,font=\bfseries,text=themecurve] at (0,2.45) {尝试 $x=63$：};
\node[good] (b7) at (2.15,2.45) {第7层\\63};
\node[count] (b6) at (3.65,2.45) {第6层\\32};
\node[count] (b5) at (5.15,2.45) {第5层\\16};
\node[count] (b4) at (6.65,2.45) {第4层\\8};
\node[count] (b3) at (8.15,2.45) {第3层\\4};
\node[count] (b2) at (9.65,2.45) {第2层\\2};
\node[count] (b1) at (11.15,2.45) {根\\1};
\draw[up] (b7)--(b6);
\draw[up] (b6)--(b5); \draw[up] (b5)--(b4); \draw[up] (b4)--(b3); \draw[up] (b3)--(b2); \draw[up] (b2)--(b1);
\node[anchor=west,font=\small,text=themecurve] at (12.05,2.45) {$\sum=126$};

\draw[themegray,line width=.6pt] (0,1.25)--(14.7,1.25);
\node[anchor=west,font=\small] at (0,.65)
  {一般地：第 $k$ 层放 $x$ 个节点，就从 $x$ 开始逐层向上取整计算最少父节点。};
\node[anchor=east,font=\bfseries,text=themecurve] at (14.7,.65)
  {本例最大为 $63$};

\end{tikzpicture}
\end{CJK*}
\end{document}
